\documentclass[10pt,a4paper]{amsart}
\usepackage[margin=19mm]{geometry}
\usepackage{amsmath,amssymb,mathtools}
\usepackage[scr=rsfs,cal=euler]{mathalfa}
\usepackage{xfrac}
\usepackage[unicode,hidelinks,bookmarks=false]{hyperref}
\newcommand{\ie}{i.e.\ }
\newcommand{\cf}{cf.\ }
\newcommand{\resp}{respectively\ }
\newcommand{\Subsection}{\subsection}
\providecommand{\nobreakdash}{\nobreak\hbox{-}}
\setlength{\emergencystretch}{2em}
\multlinegap0pt
\usepackage{amssymb}
\usepackage{mathtools}
\usepackage[shortlabels]{enumitem}
\newtheorem{proposition}{Proposition}
\newcommand{\E}{\mathbb E}
\newcommand{\tr}{\operatorname{tr}}
\title{A note on several inverse problems with generally random coefficients}
\author{C\u{a}t\u{a}lin I. C\^arstea}
\date{Source: arXiv:2605.20004v2 (CC BY 4.0)}
\begin{document}
\maketitle

\noindent\emph{Source and licence.} A note on several inverse problems with generally random coefficients. Version arXiv:2605.20004v2 (CC BY 4.0), distributed by arXiv under the Creative Commons Attribution 4.0 International licence, http://creativecommons.org/licenses/by/4.0/, as recorded in the arXiv licence metadata for this exact version and shown on the abstract page https://arxiv.org/abs/2605.20004v2. Full licence notice: \texttt{licenses/arxiv-2605.20004-CC-BY-4.0.txt}.\par\smallskip

\noindent\textit{Editorial scope.} Counted unit: the article's proposition prop:conductivity-first-invariant (source0.tex lines 1103--1134). The article's complete original proof of it is delivered with it (source0.tex lines 1136--1209, one \texttt{proof} environment of 74 lines). No mathematics is written, repaired or re-derived: the statement and the proof are the article's own lines, in the article's own notation, and no retained line is substituted or re-flowed.  No further source-local context is needed: every cross-reference of every retained line resolves inside the two delivered spans.  Nothing was omitted from any retained span, so the package carries no omission record.  No retained line contains a citation, so the wrapper carries no bibliography and no bibliography entry.  No retained line contains a figure environment, an image-inclusion command or drawing code, so no image is delivered and the package declares no asset dependency.  Editorial formatting only, in addition: an AMS \texttt{amsart} preamble that restates the article's own declarations and macros used by the retained lines, namely the environments \texttt{proposition} on separate counters, together with the standard packages those lines need.  The retained environments print in this document as Proposition 1 in order of appearance, whereas the article prints the same items with the numbering of its own full document.  The complete native source of the article as distributed in the arXiv e-print for this exact version is bundled unchanged as \texttt{sources/200961/source0.tex}.
\begin{proposition}\label{prop:conductivity-first-invariant}
Let
\[
  b(\omega,x)=\nabla\log a(\omega,x),
  \qquad
  \overline b(x)=\frac{\nabla m(x)}{m(x)},
\]
and define the positive semidefinite matrix
\begin{equation}\label{eq:conductivity-covariance-matrix}
  C_a(x)=\E\left[a(\omega,x)
  (b(\omega,x)-\overline b(x))\otimes
  (b(\omega,x)-\overline b(x))\right].
\end{equation}
The order $-4$ symbol of $\overline R-R_{\gamma_h}$ determines, for every $\theta\in S^{n-1}$,
\begin{equation}\label{eq:conductivity-invariant}
  \mathcal V_\gamma(x,\theta)
  =\tr C_a(x)-\theta^TC_a(x)\theta.
\end{equation}
Equivalently,
\[
  \mathcal V_\gamma(x,\theta)
  =\E\left[\frac{|\Pi_\theta\nabla a(\cdot,x)|^2}{a(\cdot,x)}\right]
  -\frac{|\Pi_\theta\nabla m(x)|^2}{m(x)}.
\]
More precisely,
\[
  \mathcal V_\gamma(x,\theta)
  =- |\xi|^4\sigma_{-4}(\overline R-R_{\gamma_h})(x,\xi),
  \qquad \theta=\xi/|\xi|.
\]
If $n\ge2$, knowing $\mathcal V_\gamma(x,\theta)$ for all $\theta\in S^{n-1}$ determines the whole matrix $C_a(x)$.
\end{proposition}

\begin{proof}
For the localized interior symbol we use the Kohn--Nirenberg convention
$D_x=(1/i)\partial_x$ and write the parametrix symbol as
$r\sim r_{-2}+r_{-3}+r_{-4}+\cdots$.  For a deterministic
$a=\gamma^{-1}$ the symbol of $-\nabla\cdot\gamma\nabla$ is
\[
  p=p_2+p_1,
  \qquad
  p_2=a^{-1}|\xi|^2,
  \qquad
  p_1=ia^{-2}\nabla a\cdot\xi .
\]
The equation $p\# r=1$ gives
\[
  \sum_\alpha\frac1{\alpha!}\partial_\xi^\alpha p\,D_x^\alpha r=1.
\]
At the first two orders one obtains
\[
  r_{-2}=a|\xi|^{-2},
  \qquad
  r_{-3}=i\nabla a\cdot\xi\,|\xi|^{-4}.
\]
At the next order the terms of total degree $-2$ in $p\# r$ give
\begin{equation}\label{eq:conductivity-rminus4-recursion}
\begin{split}
0={}&p_2r_{-4}+p_1r_{-3}
 +\sum_j(\partial_{\xi_j}p_2)D_{x_j}r_{-3}
 +\sum_j(\partial_{\xi_j}p_1)D_{x_j}r_{-2}  \\
&+\frac12\sum_{i,j}(\partial_{\xi_i}\partial_{\xi_j}p_2)
  D_{x_i}D_{x_j}r_{-2}.
\end{split}
\end{equation}
Substituting $r_{-2}$ and $r_{-3}$ into \eqref{eq:conductivity-rminus4-recursion} yields
\[
\begin{split}
 p_1r_{-3}&=-\frac{(\nabla a\cdot\xi)^2}{a^2|\xi|^4},\qquad
 \sum_j(\partial_{\xi_j}p_2)D_{x_j}r_{-3}
 =\frac{2\nabla^2a(\xi,\xi)}{a|\xi|^4},\\
 \sum_j(\partial_{\xi_j}p_1)D_{x_j}r_{-2}
 &=\frac{|\nabla a|^2}{a^2|\xi|^2},\qquad
 \frac12\sum_{i,j}(\partial_{\xi_i}\partial_{\xi_j}p_2)
  D_{x_i}D_{x_j}r_{-2}
 =-\frac{\Delta a}{a|\xi|^2}.
\end{split}
\]
Since $p_2^{-1}=a|\xi|^{-2}$, this gives
\begin{align*}
 \sigma_{-4}(R_\gamma)(x,\xi)=r_{-4}
 &=\frac{\Delta a}{|\xi|^4}
 -\frac{2\nabla^2a(\xi,\xi)}{|\xi|^6}
 +\frac{(\nabla a\cdot\xi)^2}{a|\xi|^6}
 -\frac{|\nabla a|^2}{a|\xi|^4}.
\end{align*}
The first two terms in $\sigma_{-4}(R_\gamma)$ are linear in $a$ and cancel when the averaged symbol is compared with the symbol obtained by replacing $a$ with $m=\E a$.  The remaining terms give
\[
\begin{split}
- |\xi|^4\sigma_{-4}(\overline R-R_{\gamma_h})(x,\xi)
&=\E\left[\frac{|\nabla a|^2}{a}
     -\frac{(\nabla a\cdot\theta)^2}{a}\right] \\
&\quad -\left(\frac{|\nabla m|^2}{m}
     -\frac{(\nabla m\cdot\theta)^2}{m}\right),
\end{split}
\]
which is the displayed expression involving $\Pi_\theta$.  Since $\nabla a=ab$ and $\nabla m=m\overline b$, this expression is also
\[
  \tr C_a(x)-\theta^TC_a(x)\theta.
\]
The matrix $C_a(x)$ is positive semidefinite by definition.  Finally, if $n\ge2$, then averaging \eqref{eq:conductivity-invariant} over $S^{n-1}$ determines $\tr C_a(x)$, since
\[
  \int_{S^{n-1}}\theta^TC_a(x)\theta\,d\theta
  =\frac1n\tr C_a(x)
\]
for normalized surface measure.  Hence $\theta^TC_a(x)\theta=\tr C_a(x)-\mathcal V_\gamma(x,\theta)$ is known for all $\theta$, and this determines $C_a(x)$.
\end{proof}

\end{document}
